\chapter{Wyznaczniki macierzy}
\label{ch:lin6-wyznaczniki}
\index{wyznacznik@Wyznacznik}
\index{permutacja@Permutacja}
\index{rozwinięcie Laplace'a@Rozwinięcie Laplace'a}

\noindent\textbf{Cel rozdziału.}
W Rozdziale~\ref{ch:lin5-eliminacja-gaussa}
nauczyliśmy się rozstrzygać odwracalność macierzy
przez eliminację. W tym rozdziale skonstruujemy
skalar, który odróżnia macierze odwracalne
od nieodwracalnych. Wyjaśnimy, dlaczego
wyznacznik zmienia znak przy zamianie wierszy,
dlaczego jest mnożony przy składaniu
odwzorowań oraz skąd pochodzą rozwinięcie
Laplace'a, macierz dołączona i wzory Cramera.
Każdy wynik udowodnimy nad dowolnym ciałem.

W całym rozdziale \(K\) jest ciałem,
\(0_K\ne1_K\), a wszystkie macierze
są macierzami nad \(K\).
Nie wymagamy, aby charakterystyka \(K\)
była różna od \(2\). Przyjmujemy
\(\mathbb N_0=\{0,1,2,\ldots\}\).
Pusty iloczyn w \(K\) ma wartość \(1_K\).
Przy \(n=0\) istnieje dokładnie jedna
macierz kwadratowa rozmiaru \(0\times0\)
i jest ona równa \(I_0\).

\section{Permutacje i ich znaki}
\label{sec:lin6-permutacje}

\begin{definicja}[Permutacja i liczba inwersji]
\label{def:lin6-permutacja-znak}
Niech \(n\in\mathbb N_0\).
\emph{Permutacją} zbioru
\(\{1,\ldots,n\}\) nazywamy jego bijekcję
na siebie. Zbiór tych permutacji
oznaczamy przez \(S_n\).
Dla \(\sigma\in S_n\) liczbą
\emph{inwersji} nazywamy liczbę
\[
 \operatorname{inv}(\sigma)=
 \#\{(i,j):1\leq i<j\leq n,\ 
                 \sigma(i)>\sigma(j)\}.
\]
\emph{Znakiem permutacji} jest liczba całkowita
\[
 \operatorname{sgn}(\sigma)
 :=(-1)^{\operatorname{inv}(\sigma)}
 \in\{-1,1\},
\]
którą we wzorach nad \(K\) rozumiemy
jako odpowiednio \(-1_K\) lub \(1_K\).
Dla \(n=0\) istnieje jedna permutacja
pustego zbioru; jej znak wynosi \(1\).
\end{definicja}

\begin{stwierdzenie}[Mnożenie znaków]
\label{prop:lin6-znak-zlozenia}
Dla \(n\in\mathbb N_0\) oraz dowolnych
\(\sigma,\tau\in S_n\) zachodzi
\[
 \operatorname{sgn}(\sigma\circ\tau)
 =\operatorname{sgn}(\sigma)
  \operatorname{sgn}(\tau),\qquad
 \operatorname{sgn}(\sigma^{-1})
 =\operatorname{sgn}(\sigma).
\]
Każda transpozycja dwóch różnych
elementów zbioru \(\{1,\ldots,n\}\)
ma znak \(-1\).
\end{stwierdzenie}
\begin{proof}
Niech \(s_k\) zamienia miejscami
\(k\) i \(k+1\), a pozostałe elementy
pozostawia bez zmian. W zapisie
\((\sigma(1),\ldots,\sigma(n))\)
złożenie \(\sigma\circ s_k\)
zamienia dwie sąsiednie wartości.
Dla par złożonych z tych dwóch pozycji
status inwersji zmienia się dokładnie raz;
dla każdej innej pozycji łączna liczba
inwersji z tymi dwiema pozostaje ta sama.
Zatem
\(\operatorname{sgn}(\sigma\circ s_k)
=-\operatorname{sgn}(\sigma)\).

Każdą permutację \(\tau\) można zapisać
jako iloczyn takich zamian: kolejnymi
zamianami sąsiednich wartości usuwamy
inwersje w zapisie \(\tau\), aż uzyskamy
permutację identycznościową. Każdy taki krok
zmniejsza liczbę inwersji dokładnie o \(1\),
stąd liczba użytych kroków ma tę samą
parzystość co \(\operatorname{inv}(\tau)\).
Zastosowanie tych samych zamian do
zapisu \(\sigma\) zmienia jego znak
o czynnik \((-1)^{\operatorname{inv}(\tau)}\),
co daje pierwszą równość.

Ponieważ znak permutacji identycznościowej
wynosi \(1\), podstawienie
\(\tau=\sigma^{-1}\) daje
\(1=\operatorname{sgn}(\sigma)
\operatorname{sgn}(\sigma^{-1})\);
obie liczby należą do \(\{-1,1\}\),
zatem są równe.
Transpozycja pozycji \(i<j\) jest
iloczynem \(2(j-i)-1\) zamian sąsiednich,
więc ma znak \(-1\).
Dla \(n=0\) lub \(n=1\) jedyną
permutacją jest identyczność i
wszystkie stwierdzenia są spełnione.
\end{proof}

\begin{przyklad}[Znaki konkretnych permutacji]
\label{prz:lin6-znaki}
Dla \(\sigma=(2,3,1)\in S_3\),
zapisanej przez kolejne wartości,
inwersjami są pary \((1,3)\) i \((2,3)\),
więc \(\operatorname{sgn}(\sigma)=1\).
Dla \(\tau=(2,1,3)\) występuje
dokładnie jedna inwersja, zatem
\(\operatorname{sgn}(\tau)=-1\).
To rozróżnienie będzie odpowiadało
znakom składników wyznacznika.
\end{przyklad}

\section{Definicja wyznacznika}
\label{sec:lin6-definicja}

\begin{definicja}[Wyznacznik macierzy kwadratowej]
\label{def:lin6-wyznacznik}
Niech \(n\in\mathbb N_0\) oraz
\(A=(a_{ij})\in\operatorname{Mat}_{n\times n}(K)\).
\emph{Wyznacznikiem} macierzy \(A\)
nazywamy skalar
\[
 \det A :=
 \sum_{\sigma\in S_n}
 \operatorname{sgn}(\sigma)
 \prod_{i=1}^n a_{i,\sigma(i)}
 \in K.
\]
Dla \(n=0\) suma zawiera jeden
składnik, odpowiadający pustej permutacji,
i przyjmujemy \(\det I_0=1_K\).
\end{definicja}

\begin{przyklad}[Wyznaczniki rozmiaru jeden i dwa]
\label{prz:lin6-male}
Dla \(a,b,c,d\in K\) mamy
\[
 \det(a)=a,\qquad
 \det\begin{pmatrix}a&b\\c&d\end{pmatrix}
 =ad-bc.
\]
Drugi wzór wynika z dwóch permutacji
zbioru \(\{1,2\}\), o znakach \(1,-1\).
W ciele charakterystyki \(2\) wyrażenie
\(-bc\) jest równe \(bc\); definicja
i wszystkie późniejsze twierdzenia
pozostają prawdziwe.
\end{przyklad}

\begin{twierdzenie}[Wieloliniowość, naprzemienność i normalizacja]
\label{thm:lin6-wlasnosci-aksjomatyczne}
Niech \(n\in\mathbb N_0\).
Wyznacznik, traktowany jako funkcja
\(n\) kolejnych wierszy
\((r_1,\ldots,r_n)\in(K^n)^n\),
jest liniowy w każdym wierszu,
przy ustaleniu pozostałych.
Jeżeli dwa różne wiersze są równe,
to wyznacznik jest równy \(0_K\).
Ponadto \(\det I_n=1_K\).
Dla \(n\geq1\) te trzy warunki
jednoznacznie wyznaczają funkcję
\(\det\colon\operatorname{Mat}_{n\times n}(K)\to K\).
\end{twierdzenie}
\begin{proof}
W każdej sumie
\(\prod_i a_{i,\sigma(i)}\)
występuje dokładnie jeden wyraz
z danego wiersza. Rozdzielność
mnożenia w \(K\) pozwala więc wyodrębnić
współczynniki kombinacji liniowej
tego wiersza, co daje liniowość.

Załóżmy, że wiersze o indeksach
\(p\ne q\) są jednakowe.
Zgrupujmy wyrazy definicji w pary
odpowiadające permutacjom
\(\sigma\) i \(\sigma\circ(p\,q)\).
Iloczyny skalarów w obu składnikach
są takie same, bo zamieniamy tylko
dwa identyczne wiersze. Znaki są
przeciwne na mocy
Stwierdzenia~\ref{prop:lin6-znak-zlozenia}.
W każdej parze składniki sumują się do
zera. Permutacje w parze są różne,
zatem parowanie jest poprawne.
Rozumowanie działa również w
charakterystyce \(2\), gdyż \(u+(-u)=0_K\).

W macierzy \(I_n\) iloczyn odpowiadający
permutacji różnej od identyczności
zawiera co najmniej jeden wyraz zerowy.
Składnik permutacji identycznościowej
wynosi \(1_K\), stąd \(\det I_n=1_K\).

Dla jednoznaczności niech \(F\) ma te
same trzy własności. Z naprzemienności
i wieloliniowości wynika zmiana znaku
przy zamianie dowolnych dwóch wierszy:
podstawienie w te miejsca \(u+v\) oraz
\(u+v\) daje zero, a po rozwinięciu
zostają dwa składniki, odpowiadające
kolejnościom \((u,v)\) i \((v,u)\).
Zatem ich wartości są przeciwne,
również w charakterystyce \(2\).

Niech \(e_1,\ldots,e_n\) będzie bazą
standardową \(K^n\). Rozkładając każdy
wiersz \(r_i=\sum_j a_{ij}e_j\)
i stosując wieloliniowość, otrzymujemy
\[
 F(r_1,\ldots,r_n)
 =\sum_{j_1,\ldots,j_n=1}^n
    \left(\prod_{i=1}^n a_{i,j_i}\right)
    F(e_{j_1},\ldots,e_{j_n}).
\]
Jeżeli dwa indeksy \(j_i\) są równe,
odpowiadający składnik znika.
Pozostałe ciągi są permutacjami
\((\sigma(1),\ldots,\sigma(n))\).
Własność zmiany znaku przy zamianie
wierszy oraz \(F(e_1,\ldots,e_n)=1_K\)
dają
\(F(e_{\sigma(1)},\ldots,e_{\sigma(n)})
=\operatorname{sgn}(\sigma)1_K\).
Wzór na \(F\) staje się dokładnie
wzorem z Definicji~\ref{def:lin6-wyznacznik}.
Dla \(n=0\) normalizacja ustala
jedyną wartość funkcji.
\end{proof}

\begin{wniosek}[Wpływ operacji na wierszach]
\label{cor:lin6-operacje}
Niech \(n\in\mathbb N_0\),
\(A\in\operatorname{Mat}_{n\times n}(K)\),
a \(B\) będzie wynikiem jednej dozwolonej
operacji elementarnej na wierszach \(A\).
Wówczas:
\begin{enumerate}[label=(\roman*)]
\item zamiana dwóch różnych wierszy daje
\(\det B=-\det A\);
\item pomnożenie jednego wiersza
przez \(\lambda\in K\setminus\{0_K\}\)
daje \(\det B=\lambda\det A\);
\item dodanie do jednego wiersza
\(\lambda\) razy innego wiersza
nie zmienia wyznacznika.
\end{enumerate}
Przy \(n=0\) operacji nie ma,
a \(\det I_0=1_K\).
\end{wniosek}
\begin{proof}
Pierwszy wzór wynika z naprzemienności
i wieloliniowości, jak w dowodzie
Twierdzenia~\ref{thm:lin6-wlasnosci-aksjomatyczne}.
Drugi to liniowość w jednym wierszu.
W trzecim przypadku wieloliniowość
daje sumę \(\det A\) i
\(\lambda\) razy wyznacznik macierzy
o dwóch identycznych wierszach;
drugi składnik jest zerowy.
\end{proof}

\begin{uwaga}[Znaki i charakterystyka dwa]
\label{rem:lin6-char2}
Równość \(\det B=-\det A\)
przy zamianie wierszy nie wystarcza
sama do udowodnienia, że dwa równe
wiersze dają zerowy wyznacznik:
w charakterystyce \(2\) mamy
\(-u=u\) dla każdego \(u\in K\).
Dlatego naprzemienność wykazaliśmy
bezpośrednio przez parowanie
przeciwnych składników.
\end{uwaga}

\section{Macierze trójkątne i transpozycja}
\label{sec:lin6-trojkatne}

\begin{twierdzenie}[Wyznacznik macierzy trójkątnej]
\label{thm:lin6-trojkatna}
Niech \(n\in\mathbb N_0\) i
\(A=(a_{ij})\in\operatorname{Mat}_{n\times n}(K)\).
Jeżeli \(A\) jest górnotrójkątna
(\(a_{ij}=0_K\) dla \(i>j\)),
to
\[
 \det A=\prod_{i=1}^n a_{ii}.
\]
Ta sama równość zachodzi dla
macierzy dolnotrójkątnej.
W szczególności \(\det I_n=1_K\).
\end{twierdzenie}
\begin{proof}
Dla macierzy górnotrójkątnej
niezerowy składnik sumy definicyjnej
wymaga \(\sigma(i)\geq i\)
dla wszystkich \(i\).
Ponieważ \(\sigma\) jest permutacją,
\(\sum_i\sigma(i)=\sum_i i\);
zatem wszystkie te nierówności
muszą być równościami.
Jedyną możliwą permutacją jest
identyczność, której znak wynosi \(1\).
Pozostaje więc wyłącznie iloczyn
wyrazów diagonalnych.

Dla macierzy dolnotrójkątnej
niezerowy składnik wymaga
\(\sigma(i)\leq i\) dla wszystkich \(i\).
Ten sam argument z sumą indeksów
znów daje permutację identycznościową.
Przy \(n=0\) wyrażenie jest
pustym iloczynem \(1_K\).
\end{proof}

\begin{stwierdzenie}[Wyznacznik transpozycji]
\label{prop:lin6-transpozycja}
Dla \(n\in\mathbb N_0\) i każdej
\(A\in\operatorname{Mat}_{n\times n}(K)\)
zachodzi
\[
 \det(A^{\mathsf T})=\det A.
\]
W szczególności wyznacznik jest
liniowy w każdej kolumnie,
znika przy dwóch identycznych kolumnach,
a zamiana, skalowanie i dodawanie
wielokrotności kolumn zmieniają go
według tych samych reguł co operacje
na wierszach.
\end{stwierdzenie}
\begin{proof}
Dla każdej permutacji \(\sigma\)
składnik sumy definiującej
\(\det(A^{\mathsf T})\) jest równy
\[
 \operatorname{sgn}(\sigma)
 \prod_{i=1}^n a_{\sigma(i),i}.
\]
Po podstawieniu \(\tau=\sigma^{-1}\)
i przestawieniu czynników
otrzymujemy
\(\operatorname{sgn}(\tau)
\prod_{j=1}^n a_{j,\tau(j)}\),
gdyż
\(\operatorname{sgn}(\tau)=
\operatorname{sgn}(\sigma)\).
Odwzorowanie \(\sigma\mapsto\sigma^{-1}\)
jest bijekcją \(S_n\), więc sumy są równe.
Własności kolumn wynikają z odpowiednich
własności wierszy po transpozycji.
\end{proof}

\begin{przyklad}[Obliczenie wyznacznika przez eliminację]
\label{prz:lin6-eliminacja}
Niech
\[
 A=\begin{pmatrix}
 1&2&0\\
 0&1&3\\
 2&0&1
 \end{pmatrix}
 \in\operatorname{Mat}_{3\times3}(\mathbb R).
\]
Operacje
\(R_3\leftarrow R_3-2R_1\)
oraz
\(R_3\leftarrow R_3+4R_2\)
dają
\[
 A\sim
 \begin{pmatrix}
 1&2&0\\
 0&1&3\\
 0&-4&1
 \end{pmatrix}
 \sim
 \begin{pmatrix}
 1&2&0\\
 0&1&3\\
 0&0&13
 \end{pmatrix}.
\]
Są to dodania wielokrotności
innych wierszy, więc nie zmieniają
wyznacznika. Zatem \(\det A=13\).
Ten sam rachunek działa nad dowolnym ciałem,
jeśli liczby całkowite interpretujemy
jako wielokrotności \(1_K\).
W charakterystyce \(13\) wyznacznik
wynosi \(0_K\), nie \(13\ne0\).
\end{przyklad}

\section{Wyznacznik iloczynu i odwracalność}
\label{sec:lin6-mnozenie}

\begin{twierdzenie}[Multiplikatywność wyznacznika]
\label{thm:lin6-multiplikatywnosc}
Dla \(n\in\mathbb N_0\) i dowolnych
\(A,B\in\operatorname{Mat}_{n\times n}(K)\)
zachodzi
\[
 \det(AB)=\det A\cdot\det B.
\]
\end{twierdzenie}
\begin{proof}
Przy ustalonej macierzy \(B\)
rozważmy funkcję
\(F(A)=\det(AB)\) zależną od
kolejnych wierszy \(A\).
Każdy wiersz \(AB\) otrzymujemy
mnożąc odpowiedni wiersz \(A\)
przez \(B\) z prawej strony;
operacja ta jest liniowa.
Zatem \(F\) jest liniowa w każdym
wierszu \(A\). Jeśli dwa wiersze \(A\)
są równe, odpowiadające wiersze \(AB\)
też są równe i \(F(A)=0_K\).

Dowód jednoznaczności funkcji
naprzemiennej i wieloliniowej z
Twierdzenia~\ref{thm:lin6-wlasnosci-aksjomatyczne}
pokazuje, że dla dowolnej takiej funkcji
\(F(A)=F(I_n)\det A\).
Tutaj \(F(I_n)=\det B\), a więc
\(\det(AB)=\det B\det A\).
Przemienność mnożenia skalarów w \(K\)
daje zapis z tezy.
Dla \(n=0\) obie strony są równe \(1_K\).
\end{proof}

\begin{twierdzenie}[Kryterium odwracalności]
\label{thm:lin6-odwracalnosc}
Niech \(n\in\mathbb N_0\) i
\(A\in\operatorname{Mat}_{n\times n}(K)\).
Następujące warunki są równoważne:
\begin{enumerate}[label=(\roman*)]
\item \(\det A\ne0_K\);
\item \(A\) jest odwracalna;
\item \(\operatorname{rk}A=n\);
\item kolumny \(A\) tworzą bazę \(K^n\).
\end{enumerate}
Jeżeli te warunki zachodzą, to
\[
 \det(A^{-1})=(\det A)^{-1}.
\]
\end{twierdzenie}
\begin{proof}
Jeżeli \(A\) jest odwracalna,
to z mnożenia wyznaczników mamy
\[
 1_K=\det I_n
    =\det(AA^{-1})
    =\det A\det(A^{-1}).
\]
Zatem \(\det A\ne0_K\) i otrzymujemy
również wzór na wyznacznik odwrotności.

Odwrotnie, załóżmy \(\det A\ne0_K\).
Sprowadźmy \(A\) przez operacje
elementarne do postaci schodkowej
zredukowanej \(R\), korzystając z
Twierdzenia~\ref{thm:lin5-gauss}.
Każda operacja zmienia wyznacznik
przez pomnożenie przez niezerowy
skalar albo pozostawia go bez zmian.
Wobec tego \(\det R\ne0_K\).
Jeżeli \(R\) miałaby mniej niż \(n\)
niezerowych wierszy, zawierałaby
wiersz zerowy i z liniowości wyznacznika
otrzymalibyśmy \(\det R=0_K\).
Stąd ma \(n\) pivotów i jest równa \(I_n\).
Wniosek~\ref{cor:lin5-pivoty-odwracalnosc}
daje odwracalność \(A\).
Równoważność odwracalności,
pełnego rzędu i istnienia bazy z kolumn
wynika z
Twierdzenia~\ref{thm:lin4-odwracalnosc-kryteria}
i wzoru wymiarowego.
Dla \(n=0\) wszystkie warunki zachodzą,
a \(I_0^{-1}=I_0\).
\end{proof}

\begin{wniosek}[Wyznacznik a podobieństwo]
\label{cor:lin6-podobienstwo}
Niech \(n\in\mathbb N_0\),
\(A\in\operatorname{Mat}_{n\times n}(K)\),
a \(P\in\operatorname{Mat}_{n\times n}(K)\)
będzie odwracalna. Wtedy
\[
 \det(P^{-1}AP)=\det A.
\]
W szczególności dla endomorfizmu
\(T\colon V\to V\) skończenie wymiarowej
przestrzeni wektorowej \(V\) nad \(K\),
wyznacznik macierzy \(T\) jest
niezależny od wyboru uporządkowanej bazy.
Liczbę tę oznaczamy przez \(\det T\).
\end{wniosek}
\begin{proof}
Z multiplikatywności wyznacznika
i poprzedniego twierdzenia mamy
\[
 \det(P^{-1}AP)
 =(\det P)^{-1}\det A\det P
 =\det A.
\]
Macierze tego samego endomorfizmu
w różnych uporządkowanych bazach
są podobne na mocy
Wniosku~\ref{cor:lin4-podobienstwo}.
Daje to niezależność \(\det T\) od bazy.
\end{proof}

\begin{uwaga}[Wyznacznik nie jest liniowy w całej macierzy]
\label{rem:lin6-nieliniowosc}
Wieloliniowość w poszczególnych wierszach
nie oznacza liniowości funkcji
\(A\mapsto\det A\). Na przykład nad
\(\mathbb R\), dla \(n=2\),
\[
 \det(I_2+I_2)=\det(2I_2)=4,
 \qquad \det I_2+\det I_2=2.
\]
Z ogólnej wieloliniowości wynika
\(\det(\lambda A)=\lambda^n\det A\),
a nie \(\lambda\det A\).
Wzór potęgowy obowiązuje też przy \(n=0\),
bo wówczas \(\lambda^0=1_K\)
(również dla \(\lambda=0_K\) przez
konwencję pustego iloczynu).
\end{uwaga}

\section{Rozwinięcie Laplace'a i macierz dołączona}
\label{sec:lin6-laplace}

\begin{definicja}[Minory i dopełnienia algebraiczne]
\label{def:lin6-minor}
Niech \(n\geq1\) i
\(A=(a_{ij})\in\operatorname{Mat}_{n\times n}(K)\).
Dla \(1\leq i,j\leq n\) przez
\(A_{\widehat i,\widehat j}\)
oznaczamy macierz rozmiaru
\((n-1)\times(n-1)\), powstałą po
usunięciu wiersza \(i\) i kolumny \(j\),
z zachowaniem kolejności pozostałych.
Liczbę
\(M_{ij}:=\det A_{\widehat i,\widehat j}\)
nazywamy \emph{minorem} odpowiadającym
pozycji \((i,j)\), a
\[
 C_{ij}:=(-1)^{i+j}M_{ij}
\]
jej \emph{dopełnieniem algebraicznym}.
Dla \(n=1\) minor jest wyznacznikiem
pustej macierzy, czyli \(1_K\).
\end{definicja}

\begin{twierdzenie}[Rozwinięcie Laplace'a]
\label{thm:lin6-laplace}
Niech \(n\geq1\) oraz
\(A\in\operatorname{Mat}_{n\times n}(K)\).
Dla każdego \(i\in\{1,\ldots,n\}\)
i każdego \(j\in\{1,\ldots,n\}\)
zachodzą wzory
\[
 \det A=\sum_{k=1}^{n} a_{ik}C_{ik},
 \qquad
 \det A=\sum_{k=1}^{n} a_{kj}C_{kj}.
\]
Są to odpowiednio rozwinięcia
według wiersza \(i\) i kolumny \(j\).
\end{twierdzenie}
\begin{proof}
W sumie definiującej wyznacznik
grupujemy permutacje według wartości
\(j=\sigma(i)\). Dla ustalonego \(j\)
przyporządkowanie \(\sigma\)
wyznacza bijekcję pomiędzy pozostałymi
\(n-1\) indeksami wierszy i kolumn,
czyli permutację indeksów macierzy
\(A_{\widehat i,\widehat j}\).

Po przeniesieniu wiersza \(i\)
i kolumny \(j\) na pierwsze miejsca
porządków indeksów wykonujemy odpowiednio
\(i-1\) i \(j-1\) przesunięć sąsiednich.
Z własności znaku permutacji
otrzymujemy
\[
 \operatorname{sgn}(\sigma)
 =(-1)^{(i-1)+(j-1)}
  \operatorname{sgn}(\widetilde\sigma)
 =(-1)^{i+j}
  \operatorname{sgn}(\widetilde\sigma),
\]
gdzie \(\widetilde\sigma\) jest
permutacją pozostałych indeksów.
Suma składników z warunkiem
\(\sigma(i)=j\) wynosi zatem
\[
 a_{ij}(-1)^{i+j}
 \det A_{\widehat i,\widehat j}
 =a_{ij}C_{ij}.
\]
Zsumowanie po \(j\) daje rozwinięcie
według wiersza \(i\).
Rozwinięcie według kolumny wynika
ze wzoru wierszowego zastosowanego
do \(A^{\mathsf T}\) oraz z równości
\(\det A^{\mathsf T}=\det A\).
Dla \(n=1\) pozostała macierz jest
pusta, więc wzór daje po prostu
\(\det(a_{11})=a_{11}\).
\end{proof}

\begin{definicja}[Macierz dołączona]
\label{def:lin6-dolaczona}
Dla \(n\geq1\) i
\(A\in\operatorname{Mat}_{n\times n}(K)\)
\emph{macierzą dołączoną} do \(A\)
nazywamy transpozycję macierzy
dopełnień algebraicznych:
\[
 \operatorname{adj}A=(C_{ji})_{1\leq i,j\leq n}.
\]
\end{definicja}

\begin{twierdzenie}[Tożsamość dla macierzy dołączonej]
\label{thm:lin6-dolaczona}
Niech \(n\geq1\) i
\(A\in\operatorname{Mat}_{n\times n}(K)\).
Wtedy
\[
 A\,\operatorname{adj}A
 =\operatorname{adj}A\,A
 =(\det A)I_n.
\]
Jeżeli \(\det A\ne0_K\), to
\[
 A^{-1}=(\det A)^{-1}\operatorname{adj}A.
\]
\end{twierdzenie}
\begin{proof}
Wyraz \((i,j)\) iloczynu
\(A\,\operatorname{adj}A\) jest
\(\sum_{k=1}^{n}a_{ik}C_{jk}\).
Dla \(i=j\) jest to rozwinięcie
Laplace'a według wiersza \(i\),
czyli \(\det A\).

Dla \(i\ne j\) zastąpmy wiersz \(j\)
macierzy \(A\) jej wierszem \(i\).
Otrzymana macierz ma dwa równe wiersze,
więc jej wyznacznik jest równy \(0_K\).
Rozwijając ten wyznacznik według
zastąpionego wiersza \(j\),
otrzymujemy dokładnie
\(\sum_k a_{ik}C_{jk}\).
Dopełnienia algebraiczne tego wiersza
są niezmienione, ponieważ powstają
po usunięciu wiersza \(j\).
Zatem elementy pozadiagonalne
iloczynu są zerowe.

Analogicznie rozwinięcie według kolumn
i zastąpienie kolumny \(i\) kolumną \(j\)
dowodzą
\(\operatorname{adj}A\,A=(\det A)I_n\).
Jeżeli wyznacznik jest niezerowy,
wolno podzielić przez \(\det A\);
otrzymujemy dwustronną odwrotność
i wzór na \(A^{-1}\).
\end{proof}

\begin{przyklad}[Macierz dołączona rozmiaru dwa]
\label{prz:lin6-adjugate-2}
Dla \(A=\begin{pmatrix}a&b\\c&d\end{pmatrix}\)
mamy
\[
 \operatorname{adj}A
 =\begin{pmatrix}d&-b\\-c&a\end{pmatrix},
 \qquad
 A\,\operatorname{adj}A=(ad-bc)I_2.
\]
Jeśli \(ad-bc\ne0_K\), to
\[
 A^{-1}=\frac{1}{ad-bc}
 \begin{pmatrix}d&-b\\-c&a\end{pmatrix}.
\]
Dla macierzy z
Przykładu~\ref{prz:lin4-odwrotnosc-2}
otrzymujemy ponownie
\(\begin{pmatrix}1&-1\\-1&2\end{pmatrix}\).
\end{przyklad}

\section{Wzory Cramera i interpretacja geometryczna}
\label{sec:lin6-cramer}

\begin{twierdzenie}[Wzory Cramera]
\label{thm:lin6-cramer}
Niech \(n\geq1\),
\(A\in\operatorname{Mat}_{n\times n}(K)\)
będzie macierzą o \(\det A\ne0_K\),
a \(b\in K^n\).
Dla \(1\leq j\leq n\) oznaczmy
przez \(A_j(b)\) macierz otrzymaną
z \(A\) po zastąpieniu \(j\)-tej
kolumny kolumną \(b\).
Układ \(Ax=b\) ma dokładnie jedno
rozwiązanie \(x\in K^n\), którego
współrzędne są równe
\[
 x_j=\frac{\det A_j(b)}{\det A},
 \qquad 1\leq j\leq n.
\]
\end{twierdzenie}
\begin{proof}
Na mocy
Twierdzenia~\ref{thm:lin6-odwracalnosc}
macierz \(A\) jest odwracalna,
więc rozwiązanie \(x=A^{-1}b\)
istnieje i jest jednoznaczne.
Niech \(a_1,\ldots,a_n\) będą
kolumnami \(A\). Równość \(Ax=b\)
oznacza
\(b=\sum_{k=1}^n x_ka_k\).

Stosując liniowość wyznacznika
w \(j\)-tej kolumnie, otrzymujemy
\[
 \begin{aligned}
 \det A_j(b)
 &=\det(a_1,\ldots,
        \textstyle\sum_{k=1}^n x_ka_k,
        \ldots,a_n)\\
 &=\sum_{k=1}^n x_k
   \det(a_1,\ldots,a_k,\ldots,a_n),
 \end{aligned}
\]
gdzie \(a_k\) w drugim wierszu
wstawiamy w miejsce \(j\)-tej kolumny.
Dla \(k\ne j\) wyznacznik
ma dwie identyczne kolumny i znika.
Dla \(k=j\) jest równy \(\det A\).
Zatem \(\det A_j(b)=x_j\det A\).
Ponieważ \(\det A\ne0_K\),
dzielenie daje wzór z tezy.
\end{proof}

\begin{przyklad}[Rozwiązanie układu wzorami Cramera]
\label{prz:lin6-cramer-2}
Niech \(K=\mathbb R\),
\[
 A=\begin{pmatrix}2&1\\1&1\end{pmatrix},
 \qquad b=\begin{pmatrix}5\\3\end{pmatrix}.
\]
Ponieważ \(\det A=2-1=1\),
układ ma jedno rozwiązanie.
Obliczamy
\[
 \det A_1(b)
 =\det\begin{pmatrix}5&1\\3&1\end{pmatrix}
 =2,\qquad
 \det A_2(b)
 =\det\begin{pmatrix}2&5\\1&3\end{pmatrix}
 =1.
\]
Zatem \(x_1=2\), \(x_2=1\).
Sprawdzenie przez podstawienie daje
\(2\cdot2+1=5\) oraz \(2+1=3\).
Wzory Cramera wymagają kwadratowej
macierzy o niezerowym wyznaczniku;
nie zastępują ogólnej eliminacji
dla układów prostokątnych lub osobliwych.
\end{przyklad}

\begin{uwaga}[Pole i orientacja nad \(\mathbb R\)]
\label{rem:lin6-geometria}
W przypadku \(\mathbb R^2\) wyznacznik
\(\det(u\ v)=u_1v_2-u_2v_1\)
jest zorientowanym polem równoległoboku
rozpiętego przez kolumny \(u,v\):
jego wartość bezwzględna jest polem,
a znak rozróżnia orientację.
Na przykład dla
\(u=(1,0)^{\mathsf T}\) i
\(v=(1,1)^{\mathsf T}\) pole wynosi \(1\).
Zamiana kolejności kolumn odwraca
orientację bez zmiany zwykłego pola.
Interpretacja pola dotyczy tu ciała
\(\mathbb R\), podczas gdy wszystkie
wcześniejsze własności algebraiczne
obowiązują dla dowolnego ciała \(K\).
\end{uwaga}

\begin{figure}[!htbp]
\centering
\begin{tikzpicture}[>=Latex,scale=1.0]
 \draw[->,gray] (-0.25,0) -- (3.6,0)
   node[below right] {\(x\)};
 \draw[->,gray] (0,-0.2) -- (0,2.35)
   node[above left] {\(y\)};
 \filldraw[fill=manifoldblue!13,draw=manifoldblue,thick]
   (0,0)--(1.45,0)--(2.9,1.45)--(1.45,1.45)--cycle;
 \draw[->,very thick,manifoldblue] (0,0)--(1.45,0)
   node[midway,below] {\(u\)};
 \draw[->,very thick,manifoldred] (0,0)--(1.45,1.45)
   node[midway,above left] {\(v\)};
 \fill (0,0) circle (1.2pt);
 \node[below left] at (0,0) {\(0\)};
\end{tikzpicture}
\caption{Równoległobok rozpięty przez
\(u=(1,0)^{\mathsf T}\) i \(v=(1,1)^{\mathsf T}\).
Macierz o tych kolumnach ma wyznacznik \(1\).
Rysunek ma charakter poglądowy.}
\label{fig:lin6-rownoleglobok}
\end{figure}

\section{Podsumowanie}
\label{sec:lin6-podsumowanie}

Wyznacznik jest jedyną znormalizowaną
funkcją wieloliniową i naprzemienną
wierszy macierzy kwadratowej.
Definicja przez permutacje działa
nad każdym ciałem, także w charakterystyce \(2\).
Operacje elementarne umożliwiają jego
obliczenie przez redukcję do macierzy
trójkątnej. Zachodzi
\(\det(AB)=\det A\det B\),
a macierz jest odwracalna dokładnie
wtedy, gdy jej wyznacznik jest niezerowy.
Rozwinięcie Laplace'a prowadzi do
macierzy dołączonej i wzoru na odwrotność,
natomiast wzory Cramera podają
współrzędne rozwiązania zgodnego układu
kwadratowego z odwracalną macierzą.
Wyznacznik endomorfizmu nie zależy
od wyboru jego uporządkowanej bazy.
